\chapter{Conclusão}
\label{cap:conclusao}

Este trabalho investigou a aplicação de dois paradigmas de redes neurais artificiais, o ADALINE Logístico e o Perceptron Multicamadas, no problema de predição de doenças cardíacas, avaliando cada arquitetura em duas versões de implementação (NumPy e scikit-learn) e sob três configurações de treinamento sistematicamente variadas. A comparação foi conduzida sobre dois conjuntos de dados com perfis clínicos distintos: o Banco 1 (UCI Heart Disease Dataset), voltado à triagem diagnóstica de doença cardíaca em pacientes de consulta ambulatorial, e o Banco 2 (Heart Failure Clinical Records Dataset), orientado ao prognóstico de mortalidade em pacientes com insuficiência cardíaca já diagnosticada. Ao todo, 24 experimentos independentes foram conduzidos, avaliados não apenas pela acurácia global, mas principalmente pelo número de Falsos Negativos, critério de maior relevância clínica em diagnóstico médico.

\section{Síntese dos Resultados}

Os resultados obtidos nos 24 experimentos revelaram padrões distintos conforme o banco de dados, a arquitetura e a configuração de treinamento, permitindo formular conclusões tanto sobre o desempenho absoluto de cada modelo quanto sobre os fatores que determinaram as diferenças observadas.

\subsection{Banco 1 --- Triagem Diagnóstica}

No Banco 1, o melhor desempenho clínico foi atingido pelo MLP NumPy na configuração C3 (3.000 épocas, $\eta = 0{,}0001$): acurácia de 87,68\% e apenas 15 Falsos Negativos no conjunto de teste, correspondendo a uma sensibilidade de 90,20\%, ou seja, 9 em cada 10 pacientes doentes foram corretamente identificados. Em paralelo, o MLP Sklearn na configuração C2 (1.000 épocas, $\eta = 0{,}001$) atingiu a maior acurácia absoluta de todo o trabalho: 88,04\%, com 20 FNs. Os dois modelos ADALINE, por sua vez, convergiram de forma estável para acurácias entre 85,87\% e 86,96\%, com FNs variando de 20 a 28, demonstrando progressão monotônica e previsível ao longo das configurações.

A diferença de apenas 1,08 pontos percentuais de acurácia entre o melhor ADALINE (NumPy C3, 86,96\%) e o melhor MLP (Sklearn C2, 88,04\%) é modesta no sentido técnico, mas a diferença de 5 FNs entre o ADALINE NumPy C3 (20 FNs) e o MLP NumPy C3 (15 FNs) é expressiva no sentido clínico: representa 5 pacientes cardíacos que deixariam de receber diagnóstico e tratamento em um conjunto de teste de 276 pacientes. No contexto da medicina preventiva, em que cada FN pode implicar uma complicação cardíaca não tratada, essa diferença não deve ser subestimada \cite{porto2019}.

\subsection{Banco 2 --- Prognóstico de Mortalidade}

No Banco 2, o padrão de resultados inverteu-se em relação ao Banco 1: os modelos MLP, que no primeiro banco tendiam a melhorar com mais épocas, degradaram-se progressivamente à medida que o treinamento se estendia. O MLP NumPy na configuração C1 (500 épocas, $\eta = 0{,}01$) ofereceu o melhor equilíbrio entre acurácia e critério clínico: 78,89\% de acurácia e 9 FNs, com sensibilidade de 68,97\%. O MLP Sklearn na configuração C3 atingiu o mínimo absoluto de FNs deste banco (8 FNs), porém ao custo de colapso de acurácia (63,33\%), tornando o MLP NumPy C1 a escolha de referência quando ambos os critérios são ponderados em conjunto. O ADALINE NumPy, mais conservador, apresentou progressão monotônica e atingiu sua melhor acurácia em C3 (75,56\%, 16 FNs), comportando-se de forma esperada e estável ao longo de todas as configurações.

O MLP Sklearn, por sua vez, demonstrou a degradação mais severa de todo o experimento: de 73,33\% de acurácia em C1 para 63,33\% em C3, configurando o pior resultado absoluto de todos os 24 experimentos. Esse resultado expõe um risco concreto do uso de otimizadores com momento acumulado em conjuntos de dados pequenos: sem volume de dados suficiente para regularizar o aprendizado, a velocidade acumulada pelo momento permite que o modelo continue ajustando seus parâmetros mesmo após a memorização dos exemplos de treino, resultando em colapso da capacidade de generalização \cite{silva2016}.

\section{Por que o ADALINE Mostrou-se Competitivo nos Dados Atuais}

Uma das conclusões mais relevantes deste trabalho é que o Perceptron Multicamadas, modelo teoricamente mais poderoso do que o ADALINE em virtude de sua capacidade de aprender representações não lineares, não produziu vantagens amplas e consistentes sobre o modelo mais simples nas condições experimentais deste estudo. Essa constatação não é um resultado negativo: ela reflete uma propriedade fundamental da modelagem preditiva, segundo a qual a escolha do modelo ideal é inseparável das características dos dados disponíveis \cite{silva2016}.

Dois fatores estruturais, confirmados empiricamente pelo experimento complementar da Seção~\ref{sec:resultados_bruto}, explicam por que o ADALINE se mostrou competitivo com o MLP neste contexto específico, limitando a vantagem prática do modelo mais complexo.

Antes de apresentá-los, cabe registrar a hipótese inicialmente cogitada e empiricamente descartada: a de que a binarização total das variáveis teria favorecido o ADALINE ao linearizar o espaço de características, reduzindo a necessidade da camada oculta do MLP. O experimento complementar (Seção~\ref{sec:resultados_bruto}) refutou essa hipótese: com dados contínuos e padronizados, o ADALINE NumPy superou o MLP NumPy no Banco~2 em todas as configurações (82,22\% contra 78,89\%--81,11\%), e o experimento revelou que a binarização havia prejudicado o próprio ADALINE de forma mais expressiva, com ganho de até 13,3 pontos percentuais no Banco~2 ao remover a binarização. A hipótese foi, portanto, descartada como fator primário.

O primeiro fator determinante é o volume insuficiente de dados para justificar a complexidade do MLP. Modelos multicamadas possuem substancialmente mais parâmetros a estimar: o MLP com 10 neurônios ocultos totaliza 371 parâmetros no Banco 1 e 161 no Banco 2, contra 36 e 15 parâmetros do ADALINE, respectivamente. Com 513 exemplos de treinamento no Banco 1, a razão entre dados e parâmetros para o MLP é de apenas 1,4, e de 1,04 no Banco 2. O ADALINE, por sua vez, apresenta razões superiores a 10 em ambos os bancos, mais de 10 vezes maior. Sob baixas razões exemplos/parâmetros, redes neurais profundas tendem a memorizar os exemplos individuais de treinamento em vez de extrair padrões generalizáveis, fenômeno denominado \textit{overfitting}, comportamento observado de forma evidente nos modelos MLP do Banco 2 e no MLP Sklearn do Banco 1 em C3 \cite{braga2007}.

O segundo fator é a instabilidade do SGD com momento em datasets pequenos. O \texttt{MLPClassifier} do scikit-learn foi configurado com \texttt{solver='sgd'} e momento padrão $\beta = 0{,}9$ \cite{pedregosa2011}. O momento acumulado, que combina o gradiente instantâneo com a velocidade das iterações anteriores, torna o otimizador suscetível a continuar realizando atualizações significativas mesmo em fases avançadas do treinamento, quando a utilidade dessas atualizações já foi esgotada. Essa característica, combinada ao volume reduzido dos conjuntos de dados, explica a queda severa do MLP Sklearn em C3 tanto no Banco 1 (acurácia recuando de 88,04\% para 82,61\%, com FNs saltando de 20 para 35) quanto no Banco 2 (acurácia caindo para 63,33\%). O gradiente descendente em lote da implementação NumPy, sem momento acumulado, convergiu de forma mais lenta e conservadora, generalizando melhor dentro das configurações testadas.

Diante desses dois fatores, conclui-se que, para os conjuntos de dados utilizados e com o pré-processamento adotado, o ADALINE Logístico constitui uma escolha competitiva, estável e tecnicamente justificada. Comparado ao MLP, oferece desempenho equivalente com menor risco de instabilidade, menor sensibilidade à escolha dos hiperparâmetros e comportamento previsível ao longo das configurações de treinamento.

\section{O Potencial do MLP e a Perspectiva de Continuidade}
\label{sec:potencial_mlp}

A constatação de que o ADALINE se mostrou competitivo nos dados atuais não implica, contudo, que o ADALINE seja a melhor escolha em cenários com condições diferentes. Pelo contrário: os próprios resultados deste trabalho fornecem evidências de que o MLP possui capacidade latente superior, que não pôde ser plenamente explorada em razão das limitações estruturais do conjunto de dados.

O Teorema da Aproximação Universal estabelece que um MLP com uma única camada intermediária de neurônios com ativação não linear é capaz de aproximar, com precisão arbitrária, qualquer função contínua em domínios compactos \cite{cybenko1989,hornik1991,braga2007}. No domínio da cardiologia clínica, isso significa que um MLP com dados suficientes poderia capturar interações complexas entre biomarcadores, como a elevação conjunta de creatinina sérica e CPK na presença de fração de ejeção reduzida, que um modelo linear como o ADALINE não consegue representar com fidelidade. Essas interações de ordem superior, que podem ser determinantes na distinção entre pacientes de alto e baixo risco, são precisamente o tipo de estrutura que as camadas ocultas do MLP são projetadas para aprender \cite{silva2016}.

O pré-processamento de binarização total, adotado neste trabalho com o intuito de auxiliar os modelos ao entregar-lhes representações previamente categorizadas segundo os parâmetros médicos das diretrizes brasileiras, representa uma fonte de limitação para o aproveitamento do potencial do MLP. Conforme demonstrado pelo experimento complementar da Seção~\ref{sec:resultados_bruto}, porém, essa não é a restrição primária: mesmo com dados contínuos, o MLP colapsou no Banco 2, evidenciando que o volume insuficiente de dados é o fator dominante. Ao converter variáveis contínuas ricas como a fração de ejeção e a creatinina sérica em indicadores binários de presença ou ausência de anormalidade, perde-se informação sobre a magnitude do desvio, informação que, em contextos clínicos, frequentemente carrega valor prognóstico significativo. Uma fração de ejeção de 49\% e uma de 20\% representam estágios completamente diferentes de disfunção sistólica, mas a representação binária as trata de forma idêntica \cite{rohde2018}.

Diante dessas considerações, a perspectiva natural de continuidade deste trabalho reside na aplicação dos mesmos modelos a conjuntos de dados mais volumosos e com representação contínua das variáveis. Nesse cenário, o comportamento esperado dos modelos seria substancialmente diferente do observado neste estudo. Com maior número de registros, a razão exemplos/parâmetros do MLP se elevaria a níveis que permitem generalização confiável; com variáveis contínuas, a camada oculta teria acesso a gradientes de informação que o espaço binário simplesmente não contém. A combinação dessas condições favoreceria o aproveitamento pleno da capacidade de representação não linear do MLP, permitindo que o ganho teórico sobre o ADALINE se materialize empiricamente de forma consistente.

Além do volume de dados, a continuidade do projeto poderia explorar estratégias de regularização não implementadas neste trabalho, como a regularização $L_2$ dos pesos e a técnica de \textit{dropout}, que desativam neurônios aleatoriamente durante o treinamento para reduzir a codependência entre unidades e mitigar o \textit{overfitting} \cite{silva2016}. A adoção de \textit{early stopping}, interrupção do treinamento com base no desempenho no conjunto de validação, também poderia conter a degradação observada no MLP Sklearn em configurações mais longas, preservando o ponto de melhor generalização antes que o modelo entre no regime de memorização. Por fim, uma busca sistemática de hiperparâmetros, incluindo o número de neurônios ocultos, a taxa de aprendizado e o algoritmo de otimização, por meio de validação cruzada estratificada permitiria identificar configurações superiores às três estudadas neste trabalho, tornando a comparação entre arquiteturas ainda mais rigorosa.

Cabe ressalvar, contudo, que os valores de épocas e taxa de aprendizado que definem as configurações C1 (500 épocas, $\eta = 0{,}01$), C2 (1.000 épocas, $\eta = 0{,}001$) e C3 (3.000 épocas, $\eta = 0{,}0001$) foram escolhidos de modo arbitrário, ancorados em faixas de valores usuais reportados na literatura e em observação empírica preliminar, sem um processo formal de otimização dos hiperparâmetros. Essa escolha, embora suficiente para os objetivos comparativos propostos, não garante que as três configurações testadas representem o melhor ponto de operação de cada modelo, podendo favorecer ou desfavorecer uma arquitetura em relação à outra de forma não intencional. Como direção futura, o emprego de meta-otimizadores baseados em modelos substitutos representa uma alternativa metodologicamente mais rigorosa para a seleção de hiperparâmetros. Nessa abordagem, denominada \textit{Sequential Model-Based Optimization} (SMBO), um modelo preditivo do desempenho do algoritmo é construído e utilizado iterativamente para guiar a busca por configurações superiores, explorando o espaço de hiperparâmetros com eficiência estatística e menor custo computacional do que estratégias exaustivas como a busca em grade \cite{trindade2019}.

A hipótese de insuficiência de dados discutida ao longo deste trabalho foi abordada predominantemente sob a perspectiva do número de pacientes disponíveis; existe, contudo, uma direção complementar de continuidade relacionada ao número de atributos por paciente. O conjunto de dados original do UCI Heart Disease Dataset disponibiliza 76 atributos brutos por registro nos arquivos de quatro instituições hospitalares (Cleveland, Hungarian, Switzerland e Long Beach VA) \cite{janosi1989}. Desse total, a literatura consolidada utiliza tradicionalmente apenas 14 atributos, conjunto que inclui o presente trabalho. Os atributos não explorados abrangem informações clínicas potencialmente relevantes, como o uso de medicamentos durante o exame de esforço (digital, betabloqueadores, nitratos e bloqueadores de canal de cálcio), o protocolo de exercício adotado, múltiplas medidas de pressão arterial durante o pico do esforço, a fração de ejeção ventricular em repouso e em exercício, e a presença de anormalidades de motilidade da parede ventricular. Como direção de continuidade, a exploração de um subconjunto mais amplo desses 76 atributos representa uma estratégia de enriquecimento informacional que não exige a coleta de novos pacientes: ao ampliar o espaço de características disponível, seria possível avaliar se o volume de informação por paciente, e não apenas o número de pacientes, constitui um fator limitante adicional ao desempenho observado, investigando hipóteses que a representação reduzida de 14 atributos não permite formular.

Em uma perspectiva mais distante do escopo deste trabalho, situa-se a possibilidade de uso de inteligência artificial de forma autônoma para a descoberta de novos fatores preditivos em cardiologia. Nesse cenário, um agente de IA seria responsável por buscar novas bases de dados clínicos disponíveis publicamente, minerar variáveis ainda não amplamente reconhecidas como preditivas pela literatura médica e configurar e avaliar automaticamente novas arquiteturas de redes neurais com o objetivo de quantificar o valor preditivo dessas variáveis. Embora tecnicamente viável com os recursos computacionais e os arcabouços de automação existentes, e cientificamente promissora pela capacidade de exploração sistemática de hipóteses em escala, essa abordagem extrapola o escopo de um trabalho de conclusão de curso, sendo mais adequada como tema de investigação em nível de pós-graduação, onde o rigor metodológico exigido para validação clínica de novos marcadores preditivos pode ser atendido com maior profundidade.

\section{Contribuições do Trabalho}

Este trabalho oferece contribuições em três dimensões complementares.

No plano técnico-educacional, a implementação dos dois modelos de forma integralmente manual, sem delegar ao scikit-learn a lógica interna de treinamento, permitiu a verificação empírica das formulações matemáticas apresentadas na fundamentação teórica. A correspondência direta entre as equações do gradiente descendente, da retropropagação e das regras de atualização de pesos e os resultados observados nas curvas de aprendizado e nas matrizes de confusão demonstra a validade das derivações apresentadas nos Capítulos~\ref{cap:fundamentacao} e~\ref{cap:metodologia}. Esse exercício de implementação do zero possui valor pedagógico intrínseco, pois expõe os mecanismos internos que as bibliotecas de alto nível abstraem.

No plano metodológico, a adoção do Falso Negativo como critério clínico diferenciador, ao lado da acurácia global, representa uma contribuição ao rigor da avaliação de modelos preditivos em contextos biomédicos. A comparação sistemática de 24 experimentos (4 modelos $\times$ 3 configurações $\times$ 2 bancos de dados) com registro completo de VP, VN, FP, FN, sensibilidade, especificidade e precisão produz um panorama experimental robusto e reprodutível, com semente aleatória fixada em todos os experimentos \cite{pedregosa2011}. O experimento complementar com dados originais sem binarização (Seção~\ref{sec:resultados_bruto}), que testou e descartou empiricamente a hipótese de que a binarização seria o fator determinante da competitividade do ADALINE, constitui também uma contribuição metodológica original: raros trabalhos comparativos entre arquiteturas incluem esse tipo de ablação sistemática do pré-processamento.

No plano clínico-prático, o pré-processamento baseado em limiares das diretrizes brasileiras de cardiologia e semiologia médica, em consulta às diretrizes mais recentes da Sociedade Brasileira de Cardiologia \cite{brandao2025,rached2025,rohde2018} e ao Tratado de Semiologia Médica \cite{porto2019}, confere ao trabalho uma ancoragem clínica que vai além do tratamento puramente estatístico dos dados. A decisão de categorizar variáveis como pressão arterial, colesterol e sódio sérico com base em critérios médicos estabelecidos, em vez de apenas em limiares estatísticos dos próprios dados, aproxima o modelo da perspectiva do especialista clínico e garante que as representações de entrada sejam clinicamente interpretáveis.

\section{Limitações do Estudo}

A interpretação dos resultados deste trabalho deve considerar as limitações metodológicas que condicionaram os experimentos realizados.

O volume reduzido dos conjuntos de dados constitui a limitação mais impactante, especialmente para o Banco 2 (299 pacientes). Com conjuntos de treinamento efetivos de 513 e 167 exemplos para os dois bancos, respectivamente, a variância dos estimadores é elevada: pequenas alterações no conjunto de treino podem resultar em diferenças significativas de desempenho, dificultando a generalização das conclusões para além das partições específicas utilizadas. A ausência de validação cruzada estratificada (\textit{k-fold cross-validation}), que repetiria o processo de particionamento com diferentes subconjuntos, impede adicionalmente a estimativa de intervalos de confiança para as métricas reportadas e aumenta o risco de que os resultados sejam específicos da partição sorteada \cite{pedregosa2011}.

A binarização total das variáveis, embora metodologicamente justificada pela clareza interpretativa e pela ancoragem clínica, limita a capacidade dos modelos de explorar a informação contínua original. Conforme verificado empiricamente no experimento complementar (Seção~\ref{sec:resultados_bruto}), essa escolha limitou a capacidade preditiva de ambos os modelos, com impacto mais pronunciado no ADALINE, cujo desempenho no Banco~2 melhorou em até 13,3 pontos percentuais com dados contínuos, embora não tenha sido o fator determinante das diferenças de desempenho relativo observadas.

Para trabalhos futuros, a incorporação de análise inferencial representa uma extensão natural dos resultados aqui apresentados. Técnicas como validação cruzada repetida (\textit{repeated k-fold}) ou reamostragem por \textit{bootstrap} permitiriam tratar os desempenhos obtidos como amostras de uma distribuição, possibilitando a aplicação de testes estatísticos como o teste de McNemar sobre as matrizes de confusão. Essa abordagem forneceria embasamento estatístico adicional às comparações realizadas, permitindo quantificar a significância das diferenças de desempenho observadas entre os modelos.

A arquitetura restrita do MLP, uma única camada oculta com exatamente 10 neurônios, não representa necessariamente o ponto ótimo do espaço de hiperparâmetros. A ausência de busca sistemática de hiperparâmetros (como validação cruzada aninhada ou \textit{random search}) impede a afirmação de que os resultados obtidos representam o desempenho máximo atingível por essa arquitetura nos dados utilizados.

Por fim, a ausência de regularização explícita nos modelos NumPy, como decaimento de pesos ($L_2$) ou \textit{dropout}, e de critério de parada antecipada (\textit{early stopping}) implica que os modelos foram treinados pelo número fixo de épocas determinado pela configuração, sem mecanismo de proteção contra \textit{overfitting} além da própria taxa de aprendizado conservadora \cite{silva2016}. A incorporação dessas estratégias em versões futuras dos modelos poderia mitigar a degradação observada nas configurações mais longas.

A inclusão da variável \texttt{time} (período de acompanhamento) no Banco~2 constitui uma potencial fonte de \textit{data leakage}. Por refletir o tempo decorrido desde o diagnóstico inicial até o desfecho (óbito ou alta), essa variável pode estar correlacionada com o resultado de forma dependente do próprio desfecho: pacientes que foram a óbito frequentemente apresentam períodos de acompanhamento mais curtos. Em um cenário real de predição preventiva, essa informação não estaria disponível ao clínico no momento da avaliação. O impacto dessa variável sobre as métricas reportadas permanece desconhecido sem um experimento de ablação específico, configurando uma limitação a ser investigada em continuidade desta pesquisa.

\section{Considerações Conclusivas}

O objetivo central deste trabalho, comparar o desempenho preditivo do ADALINE Logístico e do Perceptron Multicamadas na classificação binária de doenças cardíacas, avaliando a acurácia global e o número de Falsos Negativos como critério clínico de prioridade, foi plenamente cumprido pelos 24 experimentos conduzidos, aos quais se somam 24 experimentos adicionais com dados originais sem binarização (Seção~\ref{sec:resultados_bruto}), que confirmaram empiricamente que o volume de dados é o fator determinante das diferenças observadas. A pergunta central, se o MLP produz resultados expressivamente superiores ao ADALINE para dados clínicos tabulares de cardiopatas, obteve resposta empírica clara: nas condições experimentais deste estudo, o MLP não produziu vantagem ampla nem consistente sobre o ADALINE entre os dois bancos de dados.

O conjunto dos resultados apresentados neste trabalho permite afirmar, com respaldo empírico, que a escolha entre o ADALINE Logístico e o Perceptron Multicamadas para predição de doenças cardíacas não é determinada exclusivamente pela capacidade teórica de cada arquitetura, mas pela interação entre essa capacidade e as características concretas dos dados disponíveis. Nos conjuntos de dados utilizados, com binarização total, volumes moderados e desequilíbrio de classes, o ADALINE demonstrou ser uma escolha competitiva, estável e tecnicamente justificada, enquanto o MLP revelou seu potencial apenas parcialmente, limitado pelas condições experimentais.

A perspectiva de continuidade do projeto, com bancos de dados mais volumosos, variáveis contínuas e estratégias de regularização, aponta para um cenário em que o MLP tende a superar o ADALINE de forma consistente, materializando empiricamente a vantagem teórica que o Teorema da Aproximação Universal lhe confere \cite{braga2007}. Nesse sentido, este trabalho não encerra a investigação, mas estabelece uma base metodológica rigorosa e reprodutível sobre a qual estudos futuros poderão avançar na direção de modelos preditivos cardíacos com maior poder diagnóstico.
